\section{Transfer to inductively minimal geometries}
We use the precise convention recalled in [BC98]: a pair consisting of
a finite firm rank-$r$ geometry with connected diagram and a
flag-transitive group $G$ is called \emph{minimal} when
$|G|\leq(r+1)!$. The word refers to this numerical threshold, not to
minimality in a subgroup ordering. Inductive minimality requires that
for each connected type subgraph $J$ and every flag $F$ of complementary
type, the residue pair $(\Gamma_F,G_F)$ is minimal in this sense.
We retain the foundational classification of Buekenhout, Cara and Dehon,
in the precise form recalled by Buekenhout and Cara~[BC98]: an inductively
minimal pair of positive rank $r$ has group $\Sym(r+1)$; its connected
diagram has no induced cycle longer than three, each edge belongs to a
unique maximal clique, and each vertex belongs to one or two maximal
cliques. Every such connected diagram determines one pair up to
isomorphism. The one-vertex diagram is treated separately. The proof of
this foundational classification is not replaced here.

These diagrams are precisely line graphs of trees. To see the required
direction, form a vertex for each maximal clique. A diagram vertex in
two cliques gives an edge between their clique vertices; a vertex in
one clique gives an edge to a new leaf. Distinct cliques meet in at most
one vertex, because a common edge would violate uniqueness of its
clique. The constructed graph is simple and connected, and its line
graph is the original diagram. A shortest cycle of length at least four
would give an induced cycle on its edge labels in the diagram. A triangle
would place three pairwise adjacent labels in one clique, contradicting
unique ownership of the pairwise edges by the three cliques already
represented at the cycle vertices. The constructed graph is therefore
a tree. Its reconstruction is inverse to taking the nontrivial vertex
stars of a tree as line-graph cliques. The one-vertex diagram corresponds
to the one-edge tree.

For the geometry in Lemma~\ref{U-forest}, two remaining edges give either
a six-cycle residue or $K_{2,2}$, according as they meet or are disjoint.
Its diagram is consequently the line graph of $T$. Every residue with
connected diagram and $m$ remaining types comes from a single subtree
on $m+1$ vertices, and its group is $\Sym(m+1)$. Its order is exactly
$(m+1)!$, satisfying the stated source definition of minimality.
Thus $\Gamma(T)$ is inductively minimal.
The retained uniqueness identifies every inductively minimal pair with
one of these tree geometries.

\begin{corollary}
An inductively minimal geometry has the intersection property exactly
when its diagram is linear. In particular this holds for the
$\Sym(n)$ geometries of rank $n-1$ with $2\leq n\leq7$.
\end{corollary}
\begin{proof}
Apply the retained classification, the construction just proved, and
Theorem~\ref{U-path}. The line graph of a tree is a path exactly when the
tree itself is a path: a vertex of degree at least three would give a
triangle in the line graph.
\end{proof}

